% Generated draft: requires mathematical and source review.
\documentclass[11pt]{article}
\usepackage[margin=25mm]{geometry}
\usepackage{fontspec}
\setmainfont{Latin Modern Roman}
\usepackage{amsmath,amssymb,mathtools,unicode-math}
\setmathfont{Latin Modern Math}
\usepackage{enumitem,needspace,xurl,hyperref}
\setlength{\parindent}{0pt}
\setlength{\parskip}{5pt}
\setlength{\emergencystretch}{3em}
\newcommand{\N}{\mathbb{N}}
\newcommand{\Z}{\mathbb{Z}}
\newcommand{\Q}{\mathbb{Q}}
\newcommand{\R}{\mathbb{R}}
\newcommand{\C}{\mathbb{C}}
\newcommand{\F}{\mathbb{F}}
\newcommand{\A}{\mathbb{A}}
\newcommand{\PP}{\mathbb{P}}
\newcommand{\E}{\mathbb{E}}
\newcommand{\eps}{\varepsilon}
\newcommand{\dd}{\,\mathrm{d}}
\begin{document}
\section*{UnsolvedMath ID: 800771}
\textbf{Source problem code:} catalogue entry
\subsection{Definitions and mathematical statement}
A Hirzebruch surface is
\[
\mathbb F_m=\PP_{\PP^1}\bigl(\mathcal O_{\PP^1}\oplus\mathcal O_{\PP^1}(m)\bigr),\qquad m\in\Z_{\geq 0}.
\]
For a smooth complex projective variety \(X\), let \(\mathcal G(X)\) denote a specified class of admissible Hermitian metrics. Once the curvature tensor \(R_g\) and its sign convention have been fixed, the holomorphic sectional curvature in a nonzero direction \(v\in T_x^{1,0}X\) is
\[
H_g(x,v)=\frac{R_g(v,\overline v,v,\overline v)}{g(v,\overline v)^2}.
\]
Use the sign convention in which the Fubini--Study metric has positive holomorphic sectional curvature. If the admissible metrics are K\"ahler, \(R_g\) is the usual K\"ahler curvature tensor; for general Hermitian metrics the curvature connection must also be specified. Define
\[
\begin{aligned}
P_{\geq0}^{\mathcal G}(X)&\Longleftrightarrow
\exists g\in\mathcal G(X)\ \forall x\in X\ \forall v\in T_x^{1,0}X\setminus\{0\},\quad H_g(x,v)\geq0,\\
P_{>0}^{\mathcal G}(X)&\Longleftrightarrow
\exists g\in\mathcal G(X)\ \forall x\in X\ \forall v\in T_x^{1,0}X\setminus\{0\},\quad H_g(x,v)>0.
\end{aligned}
\]
The supplied Problem 6, attributed to Matsumura (2022), asks the following three questions.
\begin{enumerate}
\item Does the blow-up of a Hirzebruch surface admit positive holomorphic sectional curvature? For a specified blow-up \(\pi:X\to\mathbb F_m\), the metric-existence property under consideration is \(P_{>0}^{\mathcal G}(X)\). The source does not specify the center, the number of blow-ups, whether successive blow-ups are included, or whether the question is existential or universal over such choices and over \(m\). Consequently, this wording does not determine a single quantified blow-up assertion.
\item For each fixed dimension \(d\), especially \(d=2\) or \(d=3\), can one classify all varieties admitting semipositive, respectively positive, holomorphic sectional curvature? In the smooth complex projective setting with a specified metric class, the collections in question are
\[
\{X:\dim_{\C}X=d,\ P_{\geq0}^{\mathcal G}(X)\},\qquad
\{X:\dim_{\C}X=d,\ P_{>0}^{\mathcal G}(X)\}.
\]
The source does not specify the equivalence relation or form of classification, nor does its extracted statement explicitly settle smoothness or the admissible metric class.
\item Can one find an algebraic-geometric analogue of semipositive, respectively positive, holomorphic sectional curvature? No candidate algebraic condition or required relationship, such as equivalence or implication, is specified. This is therefore a research objective rather than a fully specified proposition.
\end{enumerate}
The projective setting comes from the supplied source title. The smooth framework above explains the curvature predicates; it does not resolve the source's unstated hypotheses. In particular, the choice between K\"ahler and general Hermitian metrics remains unresolved.
\subsection{Short English statement}
Matsumura asks whether blow-ups of Hirzebruch surfaces admit positive holomorphic sectional curvature, how to classify varieties of fixed dimension with semipositive or positive holomorphic sectional curvature, and whether these curvature conditions have algebraic-geometric analogues. The blow-up quantifiers, metric class, classification criteria, and meaning of an algebraic analogue require clarification.
\subsection{Sources}
\begin{itemize}
\item Local UnsolvedMath registry, permanent ID \(800771\), source code \texttt{catalogue entry}, title ``Positively Curved Projective Varieties --- Problem 6.'' The supplied record attributes the problem to Matsumura (2022) and labels it \texttt{open}; this is a supplied status label, not an independently verified mathematical or literature assessment.
\item Matsumura, \emph{Open Problems on Structure of Positively Curved Projective Varieties} (2022), Problem 6, \url{https://arxiv.org/abs/2211.09153}. The supplied provenance records access on 2026-10-16 and extraction by \texttt{arxiv\_latex}; these provenance claims have not been independently verified here.
\item The extracted statement includes the cross-reference ``cf.\ [Mat18b, Problem 5.3].'' No bibliographic entry resolving \texttt{Mat18b} is supplied, so its author details, title, and URL cannot be recovered from this record alone.
\end{itemize}
\end{document}
